\section{Principles}\label{sec:task2-principle}

\subsection{Coordinate Systems and the Pinhole Model}\label{sec:task2-coordinate-system}

The world coordinate system describes points on the calibration board. The camera coordinate system is centered at the optical center. The image coordinate system represents continuous image-plane coordinates, whereas the pixel coordinate system uses discrete image coordinates referenced to the upper-left image corner. Let
\(\mathbf{X}_{w}=(X_{w},Y_{w},Z_{w})^{\mathsf T}\) be a point in world coordinates and
\(\mathbf{X}_{c}=(X_{c},Y_{c},Z_{c})^{\mathsf T}\) the same point in camera coordinates. The extrinsic parameters are a rotation matrix \(\mathbf{R}\) and a translation vector \(\mathbf{t}\):
\begin{equation}
\mathbf{X}_{c}=\mathbf{R}\mathbf{X}_{w}+\mathbf{t}.
\label{eq:task2-extrinsic}
\end{equation}

Under the pinhole model, let \(x=X_{c}/Z_{c}\) and \(y=Y_{c}/Z_{c}\) be normalized image coordinates. The projection to pixel coordinates is
\begin{equation}
s
\begin{bmatrix}
u\\v\\1
\end{bmatrix}
=
\mathbf{K}
\begin{bmatrix}
\mathbf{R}&\mathbf{t}
\end{bmatrix}
\begin{bmatrix}
X_{w}\\Y_{w}\\Z_{w}\\1
\end{bmatrix},
\qquad s\ne0,
\label{eq:task2-projection}
\end{equation}
where \((u,v)\) are pixel coordinates and \(\mathbf{K}\) is the camera matrix:
\begin{equation}
\mathbf{K}=
\begin{bmatrix}
f_{x}&\gamma&c_{x}\\
0&f_{y}&c_{y}\\
0&0&1
\end{bmatrix}.
\label{eq:task2-intrinsic}
\end{equation}
Here \(f_x\) and \(f_y\) are focal lengths in pixels, \((c_x,c_y)\) is the principal point, and \(\gamma\) is the skew parameter. For ordinary mobile-camera images, \(\gamma\) is commonly close to zero.

\subsection{Lens Distortion Model}\label{sec:task2-distortion}

Real lenses make the observed image deviate from ideal pinhole projection. Let \(r^{2}=x^{2}+y^{2}\). A common radial distortion model is
\begin{equation}
\begin{aligned}
x_{\mathrm r}&=x\left(1+k_{1}r^{2}+k_{2}r^{4}+k_{3}r^{6}\right),\\
y_{\mathrm r}&=y\left(1+k_{1}r^{2}+k_{2}r^{4}+k_{3}r^{6}\right).
\end{aligned}
\label{eq:task2-radial-distortion}
\end{equation}
The radial distortion coefficients are \(k_1,k_2,k_3\). Tangential distortion can be written as
\begin{equation}
\begin{aligned}
x_{\mathrm d}&=x_{\mathrm r}+2p_{1}xy+p_{2}\left(r^{2}+2x^{2}\right),\\
y_{\mathrm d}&=y_{\mathrm r}+p_{1}\left(r^{2}+2y^{2}\right)+2p_{2}xy.
\end{aligned}
\label{eq:task2-tangential-distortion}
\end{equation}
The final pixel coordinates are
\begin{equation}
u=f_xx_{\mathrm d}+\gamma y_{\mathrm d}+c_x,\qquad
v=f_yy_{\mathrm d}+c_y.
\label{eq:task2-distorted-pixel}
\end{equation}
OpenCV commonly returns five coefficients in the order
\((k_1,k_2,p_1,p_2,k_3)\). If an interface returns a different length, the report must preserve the returned order and length.

\subsection{Planar Calibration and the Zhang Method}\label{sec:task2-zhang}

Set the checkerboard plane to \(Z_w=0\) and define board coordinates using the square size. The mapping from a point on the plane to its image is a homography:
\begin{equation}
s
\begin{bmatrix}
u\\v\\1
\end{bmatrix}
=
\mathbf{H}
\begin{bmatrix}
X_w\\Y_w\\1
\end{bmatrix},
\qquad
\mathbf{H}=\mathbf{K}
\begin{bmatrix}
\mathbf{r}_1&\mathbf{r}_2&\mathbf{t}
\end{bmatrix},
\label{eq:task2-homography}
\end{equation}
where \(\mathbf r_1\) and \(\mathbf r_2\) are the first two columns of the rotation matrix. Multiple views of the same checkerboard provide homography constraints. The orthogonality and unit-length constraints of the rotation columns jointly estimate the intrinsic parameters and then the pose of each view. A nonlinear optimization finally refines the intrinsic parameters, distortion coefficients, and extrinsic parameters by minimizing the discrepancy between observed and projected corners.

In the program, the checkerboard object points lie on \(Z_w=0\), and the detected image points form the corresponding two-dimensional observations. The function \texttt{cv2.calibrateCamera} performs the calibration from these correspondences. The checkerboard configuration and square size must match the physical board; otherwise the translation-vector scale is not physically meaningful.

\subsection{Reprojection Error}\label{sec:task2-reprojection}

Let \(\mathbf q_{j,i}\) be the observed pixel for corner \(i\) in image \(j\), and let
\(\widehat{\mathbf q}_{j,i}\) be the pixel obtained by projecting the estimated three-dimensional point with \texttt{cv2.projectPoints}. The pointwise reprojection error is
\begin{equation}
e_{j,i}=\left\lVert\widehat{\mathbf q}_{j,i}-\mathbf q_{j,i}\right\rVert_{2}.
\label{eq:task2-point-reprojection}
\end{equation}
If \(M\) is the total number of valid corners, the mean pointwise reprojection error and global RMS are
\begin{equation}
\overline e=\frac{1}{M}\sum_{j,i}e_{j,i},\qquad
e_{\mathrm{RMS}}=
\sqrt{\frac{1}{M}\sum_{j,i}e_{j,i}^{2}}.
\label{eq:task2-reprojection-error}
\end{equation}
For image \(j\), the per-image RMS and its arithmetic mean are
\begin{equation}
e_{j,\mathrm{RMS}}=
\sqrt{\frac{1}{M_j}\sum_i e_{j,i}^{2}},\qquad
\overline e_{\mathrm{view}}=
\frac{1}{n}\sum_{j=1}^{n}e_{j,\mathrm{RMS}},
\label{eq:task2-per-image-rms}
\end{equation}
where \(M_j\) is the number of corners in image \(j\) and \(n\) is the number of valid images. The NPZ fields \path{rms_error}, \path{mean_reprojection_error}, and \path{mean_per_image_rms} record these three distinct quantities.
